%% ma: L12-B18-0001
%% lop: 12
%% chuong: 6
%% bai: 18
%% dang: 12.18.1
%% loai: TLN
%% muc: TH
%% dap_an: "0,75"
%% nguon: "(NBV)-ĐỀ SỐ 3-MỖI NGÀY 1 ĐỀ THI 2025.pdf (D03, MỖI NGÀY 1 ĐỀ - Đề số 3), Phần III Câu 4"
%% kien_thuc: "Bài 18 (xác suất có điều kiện)"
%% trang_thai: chinh-thuc
%% ghi_chu: "Đề gốc cho bảng có ô chéo; đã viết lại bảng thành bảng thường. Đáp án gốc 0,75 đúng"
\begin{ex}
Để thử nghiệm tác dụng điều trị bệnh trào ngược dạ dày của hai loại thuốc $X$ và $Y$, người ta tiến hành thử nghiệm trên $2000$ người bệnh tình nguyện. Kết quả được cho trong bảng thống kê sau:
\begin{center}
\begin{tabular}{|l|c|c|}\hline
Kết quả & Dùng thuốc $X$ & Dùng thuốc $Y$\\\hline
Khỏi bệnh & $600$ & $800$\\\hline
Không khỏi bệnh & $200$ & $400$\\\hline
\end{tabular}
\end{center}
Chọn ngẫu nhiên một người bệnh tham gia thử nghiệm thuốc. Tính xác suất để người đó khỏi bệnh biết người đó sử dụng thuốc $X$ (kết quả viết dưới dạng số thập phân và làm tròn đến chữ số thập phân thứ hai).
\shortans{0,75}
\loigiai{Gọi $A$: ``người đó khỏi bệnh'', $B$: ``người đó dùng thuốc $X$''. Có $n(B)=600+200=800$ và $n(AB)=600$.
$P(A\mid B)=\dfrac{n(AB)}{n(B)}=\dfrac{600}{800}=0{,}75$.}
\end{ex}
